% ============================================================
% Discussion draft generated from the integrated evidence package.
% ============================================================

\subsection{Conservation supports a shared extracellular-matrix framework}
\label{subsec:discussion-main-findings}

The main result is not simply that the selected genes are ``highly conserved.''
Rather, several independent evidence types converge on a conserved molecular
framework: most candidates retain a secretion signal, SLRP/LRR-family
architecture, a continuous conserved alignment core, stable CDS-exon
organization, gene-specific phylogenetic placement, and---where SynVoy is
complete---locus-level support. This combination is consistent with long-term
maintenance of secreted extracellular-matrix proteins while allowing divergence
in terminal regions, individual repeats, expression level, and lineage-specific
gene history.

This interpretation agrees with the broader view of SLRPs as both structural
matrix organizers and regulators of receptor or growth-factor signalling
\citep{SchaeferIozzo2008,IozzoSchaefer2015}. A conserved N-terminal signal
peptide supports entry into the secretory pathway and therefore compatibility
with extracellular localization. Conserved LRR-rich architecture supports SLRP
family identity and preservation of a protein-interaction scaffold. Neither
feature proves that every ortholog binds the same partner or has the same
phenotype, but their joint conservation makes a conserved extracellular role
biologically plausible.

\subsection{Gene-specific relevance to cartilage and the growth plate}
\label{subsec:discussion-biological-relevance}

BGN and FMOD are the strongest anchors because their high juvenile and
growth-plate expression is accompanied by skeletal functional evidence. In mice,
combined Bgn/Fmod deficiency alters bone mass and osteoclastogenesis, and both
proteins interact with regulators of osteoclast formation \citep{Kram2017}.
FMOD is additionally localized around late-hypertrophic chondrocytes, the
secondary ossification centre, and the growth plate during knee development
\citep{Saamanen2001}. Their strong protein, tree, and gene-structure support is
therefore consistent with conserved roles at the interface between collagen-rich
matrix organization and skeletal remodelling. DCN adds a closely related class-I
comparison with direct collagen-fibril, cartilage-pericellular-matrix, and
mechanotransduction evidence. Its inclusion was methodologically important
because reciprocal comparison against DCN exposed several historical BGN
misassignments. Conserved DCN domains, alignment, tree placement, coding
organization, phenotype annotations, and the completed updated SynVoy run
support its inclusion. The completed SignalP screen detected a classical signal
peptide in 14 of 15 DCN proteins. The only negative, opossum
\texttt{XP\_001363160.3}, is a C-terminal fragment missing the corresponding
N-terminal region and therefore does not argue against secretion of intact DCN.

The FMOD locus results also illustrate a biologically meaningful consequence
of teleost duplication. Zebrafish \textit{fmodb} provided the more complete
human-FMOD sequence match, whereas \textit{fmoda} retained the ancestral
FMOD--PRELP neighbourhood. Spotted-gar \textit{fmoda} was likewise supported by
an alternative retained SynVoy locus and eight shared human flanking genes,
even though the run summary ranked a nearby lumican-like/PRELP interval first.
These cases support retention of duplicated co-ortholog evidence rather than
forcing every species into a single best-hit model. The elephant-shark FMOD-
like locus remains tentative because its reconstructed 684-aa product contains
two SLRP-like halves. Protein BLAST assigned the N-terminal half most strongly
to DCN and the C-terminal half most strongly to FMOD; two LRR N-terminal caps
and a second hydrophobic segment near residue 330 further support a compound or
fused gene model. The locus is informative, but this protein was therefore not
included in the confident FMOD sequence panel.

PRELP illustrates why candidate selection should not be reduced to one RNA
ranking. Its abundance in the broad P0 lineage sample was moderate, whereas the
direct mouse and rat growth-plate data were strong. PRELP is a cartilage matrix
protein with conserved LRR and cysteine-containing regions
\citep{Grover1996,Grover2002}; its N-terminal region and full-length protein have
also been linked to osteoclast and osteoblast regulation
\citep{Rucci2009,Li2016PRELP}. The highly conserved protein and compact
two-CDS-exon organization found here give molecular context to those reported
skeletal functions.

EPYC provides a biologically specific class-III example. Its expression was not
among the highest in the initial screen, but developmental data place EPYC in
the extracellular matrix surrounding resting, proliferating, and hypertrophic
growth-plate chondrocytes \citep{Johnson1999}. Thus, the conserved EPYC protein
architecture and phylogenetic assignment are relevant to a protein with direct
growth-plate localization. The whale-shark sequence was retained as tentative
because its conserved middle-to-C-terminal region supports EPYC identity while
earlier mismatches, reduced coverage, and an alternative reciprocal hit leave
model uncertainty.

LUM combines consistent rodent expression with a technically clean vertebrate
protein set. Experimental studies connect lumican to skeletal matrix development
and cartilage collagen deposition and fibrillogenesis
\citep{Raouf2002,Kafienah2008}. The conserved signal peptide and LRR-rich core
therefore fit a secreted collagen-associated ECM role. Nevertheless, the
lamprey placement shows that close class-II SLRP paralogs can be difficult to
separate in deep lineages. The completed LUM SynVoy run recovered best HIGH
calls in 10 of 14 targets and MEDIUM calls in four. Manual review accepted the
canonical catshark, coelacanth, and spotted-gar LUM loci because their ordered
flanking blocks agreed with reciprocal protein and tree evidence, even though
the automated summary selected different fragments. Amphioxus remained
ambiguous because its candidate interval was intergenic and lacked an
independent canonical protein. Thus, manual review refined the automated result
rather than converting all 14 targets into proven one-to-one orthologs.

OGN is scientifically useful precisely because its evidence is less uniform.
OGN processing affects its regulation of type-I collagen fibrillogenesis
\citep{Ge2004}, while vertebrate analyses document teleost duplication and
subfunctionalization \citep{Costa2018}. Those observations explain why a
divergent zebrafish protein can still be a credible OGN-family ortholog, but they
do not validate the amphioxus sequence. Excluding amphioxus from the confident
set was justified by the combination of divergent terminal structure, large
sequence-specific insertions, an alternative reciprocal hit, poor synteny, and
its position outside the former main OGN split. Spotted-gar OGN remains a more
limited uncertainty: its LRR-rich alignment is plausible, but the absence of a
SignalP-positive N terminus prevents a fully confident secretion claim.

OMD should be discussed as supporting context rather than silently promoted to
the same evidential level as the seven-gene panel. OMD is secreted by osteoblasts,
binds mineral, supports osteoblast attachment, and occurs in fetal
growth-plate-associated primary spongiosa \citep{Sommarin1998}. It has also been
linked to BMP2/SMAD osteogenic signalling through terminal LRR interactions
\citep{Lin2021OMD}. These findings make OMD relevant to the cartilage--bone
interface, but its full cross-species protein, phylogenetic, and synteny pipeline
was not performed here.

\subsection{Why comparative evolutionary analysis remains informative}
\label{subsec:discussion-evolution-orthology}

Strong conservation does not make the comparative analysis uninformative. It
establishes which molecular features are maintained, exposes lineage-specific
exceptions, and identifies annotation or orthology errors that a single BLAST
hit would miss. In this dataset, the analysis resolved an ASPN-like chicken
BGN-locus product, removed an uncertain amphioxus OGN sequence, retained a
divergent but structurally coherent zebrafish OGN, and isolated deep-lineage
cautions for spotted-gar OGN and lamprey LUM. These exceptions are more
biologically informative than a blanket statement that all sequences are
conserved.

The deep-lineage sampling also places these gene-level observations in the
published history of SLRP evolution. Tunicate SLRP-like genes provide
literature-only precursor context, while the diversification reported in
jawless vertebrates and the recognizable repertoire in cartilaginous fishes
support expansion of the family during early vertebrate evolution
\citep{Park2008,Costa2018}. The present data do not reconstruct ancestral
sequences or date individual duplications. They instead test whether modern
gene-specific assignments remain recognizable across this evolutionary
bracket. The uncertain amphioxus OGN-like candidate, tentative lamprey LUM, and
partial shark BGN models define the empirical limit of that assignment.

The species-level measurements further show why evolutionary distance and
annotation quality must be separated. Median protein identity generally rose
from jawless/cartilaginous fishes toward mammals, but spotted gar exceeded both
zebrafish and coelacanth in the descriptive median and retained five of seven
reviewed loci. Conversely, opossum proteins were relatively close to human
while every opossum synteny row was unusable because of an input-version
mismatch. Thus, low confidence in a species cannot automatically be attributed
to biological divergence, and high protein identity cannot rescue invalid
locus data. The gar--zebrafish comparison is particularly informative because
gar brackets the ray-finned vertebrate condition before teleost-specific
duplication, whereas the two zebrafish FMOD co-orthologs expose that later
duplication history.

The evidence types answer different questions. Domain architecture supports
SLRP-family membership, phylogenetic clustering supports gene-specific
assignment, conserved coding-exon organization supports gene-model plausibility,
and synteny supports locus-level continuity. Agreement among them increases
orthology confidence; disagreement identifies cases for manual review. A
two-dimensional ProtSpace projection was used only as exploratory quality
control because visual proximity in an embedding cannot by itself establish
orthology.

The extended gene-structure analysis strengthens this conclusion beyond exon
count. Conservation of intron phase at all 26 comparable coding junctions means
that splice boundaries preserve their relationship to the reading frame even
where the exact boundary moved by several residues. Together with complete,
stop-free reconstructed translations, this is strong evidence for stable coding
models and homologous gene organization. It does not imply that UTRs or isoform
usage are conserved: those features varied more strongly and depend on
annotation and tissue-specific transcript use.

The supplementary pairwise $d_N/d_S$ analysis provides a complementary
nucleotide-level view. All finite estimates were below one, which is compatible
with predominant purifying constraint and fits the conserved protein cores.
However, the many unavailable human--zebrafish estimates and the high-$d_S$
bird comparisons show why this result cannot be generalized into a claim of
uniform selection across all vertebrate branches. A branch-aware codon model on
a carefully curated ortholog set would be required for formal lineage- or
site-specific selection inference.

\subsection{What the supplementary motif and phenotype analyses add}
\label{subsec:discussion-supplementary-evidence}

The mature-protein motif analysis reinforces, but does not replace, the MSA and
phylogenetic results. Every canonical protein was most compatible with the motif
model learned from its assigned gene, including the manually reviewed tentative
records. This is useful evidence that those sequences preserve gene-associated
patterns across the mature protein. However, many recovered patterns overlap
the expected LRR- and cysteine-rich architecture. Without biochemical testing,
they should not be assigned a specific collagen-binding, receptor-binding, or
signalling function.

Curated mutant and disease annotations add an independent bridge from molecular
conservation to organismal biology. The concentration of skeletal, cartilage,
or joint terms for BGN, DCN, FMOD, and EPYC is consistent with their main-panel roles.
The broader corneal, skin, tendon, and collagen phenotypes associated with LUM
and OGN support general ECM functions but do not demonstrate a growth-plate
phenotype. PRELP's lack of a skeletal-category MGI term in the downloaded
release does not contradict its direct growth-plate expression and experimental
literature. Annotation counts are strongly influenced by study effort and
database curation and therefore cannot be compared as gene-effect sizes.

The promoter analysis produced a biologically useful negative result. Despite
testing conserved de-novo patterns and a targeted cartilage/osteogenesis panel,
no targeted motif passed correction in this small five-species-per-gene design.
Promoter sequence alone therefore does not support a shared regulatory mechanism
for the seven genes. This does not weaken the protein-conservation conclusion:
regulatory conservation can occur in distal enhancers, in cell-state-specific
chromatin, or through turnover of individual motif instances. The uncorrected
recurrences are hypotheses that would require ATAC-seq, ChIP-seq, reporter
assays, or equivalent functional evidence.

\subsection{Limitations}
\label{subsec:discussion-limitations}

The initial P0 expression samples were developmentally relevant but not
microdissected growth-plate zones, while the direct GSE114919 dataset used a
different processed scale and a limited set of rodent ages and zones. Expression
therefore supports biological relevance, not gene function or causal importance.
Cross-species equivalence of expression was not tested by pooling these scales.

The comparative panel used 13--16 representatives per gene and
five reference species for the representative-transcript gene-structure screen.
Predicted proteins, incomplete assemblies, alternative isoforms, and errors in
gene naming can all affect candidate recovery. Conserved LRR domains also occur
across related SLRP paralogs, so domain similarity is not gene-specific proof.
The extended isoform inventory is annotation-based and does not identify the
dominant transcript in growth-plate cells. Derived UTR lengths are especially
sensitive to transcript-end annotation. Pairwise NG86 values average across a
whole sequence, and $d_S$ saturation made most human--zebrafish comparisons
uninterpretable; they are not evidence for branch- or residue-specific
selection.
The compact trees are unrooted; their split support does not establish the
direction or timing of duplication events. The completed SynVoy review showed
that rescue calls and genomes with no retained HIGH or MEDIUM candidate cannot
be interpreted without annotation and protein evidence. In addition, the
opossum FASTA/GFF sequence-version mismatch made all seven current opossum
synteny rows technically uninterpretable; these require a matched-input rerun.
Finally, computational conservation
supports compatibility with known functions but cannot demonstrate collagen
binding, signalling activity, or a growth-plate phenotype in each species.

The protein-motif models were learned and evaluated on the same compact
candidate panel, and the gene-level STREME comparisons had no independent
holdout set. The promoter analysis contained only five orthologous promoters per
gene, depended on one annotated transcription start site per locus, and did not
include distal enhancers or chromatin accessibility. MGI and HPO reflect
available genotype experiments and clinical curation rather than a standardized
phenotyping experiment across all seven genes. These supplementary layers are
therefore confirmatory or hypothesis-generating rather than decisive.

\subsection{Future directions}
\label{subsec:discussion-future-directions}

SignalP coverage and focused review of the partial BGN/DCN models are complete.
All 98 current SynVoy locus decisions are recorded: 62 are accepted, 20
tentative, 12 ambiguous and four rejected. The required synteny follow-up is a matched-input
opossum rerun; the interrupted updated-tool FMOD refresh and subsequent PRELP/
EPYC comparisons are optional sensitivity analyses if they remain in the final
scope. The tentative whale-shark EPYC sequence, duplicated zebrafish FMOD calls,
and tentative elephant-shark FMOD-like model should remain visible when the
evidence table is frozen. Broader expression work would be most valuable if it added matched
juvenile growth-plate zones or well-annotated single-cell data rather than simply
adding unrelated adult tissues. Regulatory follow-up should prioritize
growth-plate-relevant ATAC-seq/ChIP-seq or reporter data rather than another
larger sequence-only motif scan. Experimental protein work could test whether
the most conserved candidates retain collagen-, mineral-, or signalling-related
activities suggested by the literature.
